\label{app:necesidades_identificadas}

Este compendio organiza las necesidades identificadas en una estructura jerárquica de tres niveles, ordenados de forma descendente conforme al nivel de detalle, desde las necesidades más generales hasta las más específicas. Al eliminar los enunciados redundantes, las 690 menciones individuales se consolidaron por similitud conceptual en 122 necesidades estandarizadas, seleccionando en cada caso la redacción más representativa. Las Tablas~\ref{tab:np01} a~\ref{tab:np14} sintetizan las 14 necesidades primarias resultantes, donde la frecuencia de cada nivel corresponde al total acumulado de las menciones que agrupa.

\Needspace{14\baselineskip}
\begin{longtable}{@{}>{\raggedright\arraybackslash}p{0.85\textwidth} r@{}}
\caption{Necesidad primaria: \emph{El robot manipula artículos domésticos.}}\label{tab:np01}\\
\toprule
\textbf{Necesidad} & \textbf{Frec.} \\
\midrule
\endfirsthead
\caption[]{Necesidad primaria: \emph{El robot manipula artículos domésticos.} (continuación).}\\
\toprule
\textbf{Necesidad} & \textbf{Frec.} \\
\midrule
\endhead
\midrule
\multicolumn{2}{r}{\small\emph{Continúa en la siguiente página}}\\
\endfoot
\bottomrule
\endlastfoot
\rowcolor{nprim} \textbf{El robot manipula artículos domésticos.} & \textbf{245} \\
\rowcolor{nsec} \textbf{El robot manipula artículos domésticos.} & \textbf{124} \\
El robot entrega artículos domésticos. & 6 \\
El robot manipula artículos ubicados por encima de su plano medio. & 18 \\
El robot manipula artículos domésticos. & 15 \\
El robot manipula artículos ubicados en zonas inaccesibles para el usuario. & 14 \\
El robot manipula artículos ubicados por debajo de su plano medio. & 12 \\
El robot manipula artículos con exactitud. & 11 \\
El robot manipula artículos pequeños. & 10 \\
El robot manipula artículos pesados. & 9 \\
El robot manipula artículos delgados. & 7 \\
El robot manipula artículos con formas diversas. & 6 \\
El robot manipula artículos grandes. & 6 \\
El robot manipula varios artículos de manera simultánea. & 3 \\
El robot manipula artículos en diferentes orientaciones. & 3 \\
El robot manipula artículos deformables. & 2 \\
El robot manipula artículos con superficies lisas. & 2 \\
\rowcolor{nsec} \textbf{El robot planea trayectorias del manipulador.} & \textbf{70} \\
El robot dispone sus manipuladores en poses compatibles con otros sistemas. & 13 \\
El robot ajusta la estrategia de manipulación de los artículos conforme al entorno. & 10 \\
El robot ajusta de forma dinámica la trayectoria de sus manipuladores conforme a la posición de los artículos. & 8 \\
El robot ajusta la estrategia de manipulación conforme a las características de los artículos. & 8 \\
El robot ajusta la velocidad de manipulación para optimizar la ejecución de cada tarea. & 6 \\
El robot coordina sus manipuladores para realizar tareas colaborativas. & 6 \\
El robot sigue trayectorias suaves y controladas con sus manipuladores. & 5 \\
El robot permite definir y utilizar poses preestablecidas para sus manipuladores. & 4 \\
El robot planifica trayectorias óptimas para manipular artículos. & 2 \\
El robot evita trayectorias de riesgo para el usuario. & 8 \\
\rowcolor{nsec} \textbf{El robot extiende el alcance efectivo de sus manipuladores.} & \textbf{21} \\
El robot extiende el alcance efectivo de sus manipuladores. & 8 \\
El robot dispone sus manipuladores en un amplio rango de poses. & 5 \\
El robot mantiene el equilibrio en todas sus poses. & 8 \\
\rowcolor{nsec} \textbf{El robot toma artículos con la fuerza necesaria.} & \textbf{24} \\
El robot toma artículos con firmeza. & 15 \\
El robot toma artículos con suavidad. & 9 \\
\rowcolor{nsec} \textbf{El robot manipula equipo para realizar tareas domésticas.} & \textbf{6} \\
El robot interactúa con puertas y cajones. & 4 \\
El robot manipula equipo para realizar tareas domésticas. & 2 \\
\end{longtable}

\Needspace{14\baselineskip}
\begin{longtable}{@{}>{\raggedright\arraybackslash}p{0.85\textwidth} r@{}}
\caption{Necesidad primaria: \emph{El robot atiende las órdenes del usuario.}}\label{tab:np02}\\
\toprule
\textbf{Necesidad} & \textbf{Frec.} \\
\midrule
\endfirsthead
\caption[]{Necesidad primaria: \emph{El robot atiende las órdenes del usuario.} (continuación).}\\
\toprule
\textbf{Necesidad} & \textbf{Frec.} \\
\midrule
\endhead
\midrule
\multicolumn{2}{r}{\small\emph{Continúa en la siguiente página}}\\
\endfoot
\bottomrule
\endlastfoot
\rowcolor{nprim} \textbf{El robot atiende las órdenes del usuario.} & \textbf{75} \\
\rowcolor{nsec} \textbf{El robot facilita la interacción con el usuario.} & \textbf{42} \\
El robot facilita la interacción con el usuario. & 14 \\
El robot complementa las indicaciones verbales con apoyos visuales. & 8 \\
El robot guía al usuario mediante indicaciones verbales. & 8 \\
El robot emite sonido. & 1 \\
El robot mantiene un estilo de comunicación uniforme. & 1 \\
El robot reporta el progreso y resultado de sus tareas. & 4 \\
El robot enfoca la atención del usuario. & 5 \\
El robot opera de forma silenciosa. & 1 \\
\rowcolor{nsec} \textbf{El robot atiende las órdenes del usuario.} & \textbf{12} \\
El robot atiende las órdenes del usuario. & 3 \\
El robot interpreta comandos verbales. & 4 \\
El robot interpreta comandos visuales. & 3 \\
El robot interpreta comandos físicos. & 2 \\
\rowcolor{nsec} \textbf{El robot genera confianza en el usuario.} & \textbf{20} \\
El robot genera confianza en el usuario. & 7 \\
El robot presenta características antropomórficas. & 13 \\
\rowcolor{nsec} \textbf{El robot reproduce música indicada por el usuario.} & \textbf{1} \\
El robot reproduce música indicada por el usuario. & 1 \\
\end{longtable}

\Needspace{14\baselineskip}
\begin{longtable}{@{}>{\raggedright\arraybackslash}p{0.85\textwidth} r@{}}
\caption{Necesidad primaria: \emph{El robot realiza tareas domésticas.}}\label{tab:np03}\\
\toprule
\textbf{Necesidad} & \textbf{Frec.} \\
\midrule
\endfirsthead
\caption[]{Necesidad primaria: \emph{El robot realiza tareas domésticas.} (continuación).}\\
\toprule
\textbf{Necesidad} & \textbf{Frec.} \\
\midrule
\endhead
\midrule
\multicolumn{2}{r}{\small\emph{Continúa en la siguiente página}}\\
\endfoot
\bottomrule
\endlastfoot
\rowcolor{nprim} \textbf{El robot realiza tareas domésticas.} & \textbf{37} \\
\rowcolor{nsec} \textbf{El robot limpia entornos domésticos.} & \textbf{25} \\
El robot limpia entornos domésticos. & 6 \\
El robot ordena artículos domésticos fuera de lugar. & 3 \\
El robot cuida del jardín. & 8 \\
El robot lava utensilios domésticos. & 4 \\
El robot aspira pisos. & 3 \\
El robot talla muebles tapizados. & 1 \\
\rowcolor{nsec} \textbf{El robot realiza tareas domésticas.} & \textbf{5} \\
El robot pinta muebles y paredes. & 2 \\
El robot cuida de las mascotas. & 1 \\
El robot lava ropa. & 2 \\
\rowcolor{nsec} \textbf{El robot prepara alimentos.} & \textbf{6} \\
El robot sirve alimentos y bebidas. & 4 \\
El robot prepara alimentos. & 2 \\
\rowcolor{nsec} \textbf{El robot asiste al usuario mientras nada.} & \textbf{1} \\
El robot asiste al usuario mientras nada. & 1 \\
\end{longtable}

\Needspace{14\baselineskip}
\begin{longtable}{@{}>{\raggedright\arraybackslash}p{0.85\textwidth} r@{}}
\caption{Necesidad primaria: \emph{El robot es robusto.}}\label{tab:np04}\\
\toprule
\textbf{Necesidad} & \textbf{Frec.} \\
\midrule
\endfirsthead
\caption[]{Necesidad primaria: \emph{El robot es robusto.} (continuación).}\\
\toprule
\textbf{Necesidad} & \textbf{Frec.} \\
\midrule
\endhead
\midrule
\multicolumn{2}{r}{\small\emph{Continúa en la siguiente página}}\\
\endfoot
\bottomrule
\endlastfoot
\rowcolor{nprim} \textbf{El robot es robusto.} & \textbf{99} \\
\rowcolor{nsec} \textbf{El robot protege sus componentes.} & \textbf{47} \\
El robot protege sus componentes contra agua. & 12 \\
El robot protege sus componentes contra impactos. & 12 \\
El robot protege sus conexiones eléctricas. & 8 \\
El robot protege sus conexiones mecánicas. & 5 \\
El robot protege sus componentes contra altas temperaturas. & 2 \\
El robot previene desconexiones mientras opera. & 2 \\
El robot soporta las cargas previstas. & 6 \\
\rowcolor{nsec} \textbf{El robot supervisa el estado de sus sistemas.} & \textbf{52} \\
El robot identifica y reporta fallas en sus sistemas. & 11 \\
El robot interrumpe sus sistemas ante situaciones de riesgo. & 10 \\
El robot supervisa el estado de sus sistemas. & 10 \\
El robot restablece sus sistemas ante fallas. & 6 \\
El robot registra la tasa de éxito en la manipulación de artículos. & 3 \\
El robot verifica que los artículos hayan sido tomados con éxito. & 12 \\
\end{longtable}

\Needspace{14\baselineskip}
\begin{longtable}{@{}>{\raggedright\arraybackslash}p{0.85\textwidth} r@{}}
\caption{Necesidad primaria: \emph{El robot tiene un diseño intuitivo.}}\label{tab:np05}\\
\toprule
\textbf{Necesidad} & \textbf{Frec.} \\
\midrule
\endfirsthead
\caption[]{Necesidad primaria: \emph{El robot tiene un diseño intuitivo.} (continuación).}\\
\toprule
\textbf{Necesidad} & \textbf{Frec.} \\
\midrule
\endhead
\midrule
\multicolumn{2}{r}{\small\emph{Continúa en la siguiente página}}\\
\endfoot
\bottomrule
\endlastfoot
\rowcolor{nprim} \textbf{El robot tiene un diseño intuitivo.} & \textbf{53} \\
\rowcolor{nsec} \textbf{El robot facilita su operación.} & \textbf{25} \\
El robot facilita su operación. & 8 \\
El robot facilita el inicio de sus sistemas. & 7 \\
El robot integra sus sistemas de forma coherente. & 4 \\
El robot facilita su configuración. & 2 \\
El robot integra todos sus componentes de hardware. & 2 \\
El robot facilita el montaje de sus componentes. & 1 \\
El robot controla su hardware con el stack de software existente. & 1 \\
\rowcolor{nsec} \textbf{El robot facilita su mantenimiento.} & \textbf{28} \\
El robot facilita su mantenimiento. & 6 \\
El robot mantiene sus componentes ordenados e identificados. & 5 \\
El robot sigue un esquema claro de diseño. & 4 \\
El robot facilita el desmontaje de sus componentes. & 3 \\
El robot opera sus componentes de forma independiente. & 7 \\
El robot incluye documentación técnica completa. & 3 \\
\end{longtable}

\Needspace{14\baselineskip}
\begin{longtable}{@{}>{\raggedright\arraybackslash}p{0.85\textwidth} r@{}}
\caption{Necesidad primaria: \emph{El robot identifica formas en el entorno.}}\label{tab:np06}\\
\toprule
\textbf{Necesidad} & \textbf{Frec.} \\
\midrule
\endfirsthead
\caption[]{Necesidad primaria: \emph{El robot identifica formas en el entorno.} (continuación).}\\
\toprule
\textbf{Necesidad} & \textbf{Frec.} \\
\midrule
\endhead
\midrule
\multicolumn{2}{r}{\small\emph{Continúa en la siguiente página}}\\
\endfoot
\bottomrule
\endlastfoot
\rowcolor{nprim} \textbf{El robot identifica formas en el entorno.} & \textbf{50} \\
\rowcolor{nsec} \textbf{El robot identifica artículos en el entorno.} & \textbf{44} \\
El robot identifica artículos en el entorno. & 25 \\
El robot determina la posición exacta de los artículos en el entorno. & 8 \\
El robot identifica artículos transparentes. & 2 \\
El robot construye modelos digitales de los artículos que identifica. & 1 \\
El robot distingue los artículos de la superficie que los soporta. & 1 \\
El robot identifica zonas de manipulación óptimas en los artículos. & 3 \\
El robot distingue superficies adyacentes. & 1 \\
El robot captura imágenes nítidas del entorno. & 3 \\
\rowcolor{nsec} \textbf{El robot identifica rostros en el entorno.} & \textbf{6} \\
El robot identifica rostros en el entorno. & 5 \\
El robot advierte al usuario de visitantes desconocidos. & 1 \\
\end{longtable}

\Needspace{14\baselineskip}
\begin{longtable}{@{}>{\raggedright\arraybackslash}p{0.85\textwidth} r@{}}
\caption{Necesidad primaria: \emph{El robot es autónomo.}}\label{tab:np07}\\
\toprule
\textbf{Necesidad} & \textbf{Frec.} \\
\midrule
\endfirsthead
\caption[]{Necesidad primaria: \emph{El robot es autónomo.} (continuación).}\\
\toprule
\textbf{Necesidad} & \textbf{Frec.} \\
\midrule
\endhead
\midrule
\multicolumn{2}{r}{\small\emph{Continúa en la siguiente página}}\\
\endfoot
\bottomrule
\endlastfoot
\rowcolor{nprim} \textbf{El robot es autónomo.} & \textbf{19} \\
\rowcolor{nsec} \textbf{El robot planea tareas.} & \textbf{12} \\
El robot modifica su plan de forma dinámica. & 7 \\
El robot genera un plan conforme a una meta establecida. & 2 \\
El robot agenda la ejecución de sus tareas. & 2 \\
El robot ayuda al usuario a organizar sus actividades. & 1 \\
\rowcolor{nsec} \textbf{El robot opera sin intervención del usuario.} & \textbf{7} \\
El robot sigue un plan sin órdenes explícitas. & 2 \\
El robot determina su posición en el entorno sin intervención del usuario. & 1 \\
El robot brinda asistencia en emergencias domésticas. & 4 \\
\end{longtable}

\Needspace{13\baselineskip}
\begin{longtable}{@{}>{\raggedright\arraybackslash}p{0.85\textwidth} r@{}}
\caption{Necesidad primaria: \emph{El robot es portátil.}}\label{tab:np08}\\
\toprule
\textbf{Necesidad} & \textbf{Frec.} \\
\midrule
\endfirsthead
\caption[]{Necesidad primaria: \emph{El robot es portátil.} (continuación).}\\
\toprule
\textbf{Necesidad} & \textbf{Frec.} \\
\midrule
\endhead
\midrule
\multicolumn{2}{r}{\small\emph{Continúa en la siguiente página}}\\
\endfoot
\bottomrule
\endlastfoot
\rowcolor{nprim} \textbf{El robot es portátil.} & \textbf{22} \\
\rowcolor{nsec} \textbf{El robot facilita su transporte.} & \textbf{18} \\
El robot es ligero. & 7 \\
El robot facilita su transporte como equipaje personal. & 6 \\
El robot facilita el transporte individual de sus componentes. & 3 \\
El robot facilita su envío por servicios de paquetería. & 2 \\
\rowcolor{nsec} \textbf{El robot tiene dimensiones compatibles con entornos domésticos.} & \textbf{4} \\
El robot tiene dimensiones compatibles con entornos domésticos. & 3 \\
El robot facilita su almacenamiento. & 1 \\
\end{longtable}

\Needspace{10\baselineskip}
\begin{longtable}{@{}>{\raggedright\arraybackslash}p{0.85\textwidth} r@{}}
\caption{Necesidad primaria: \emph{El robot navega en entornos domésticos.}}\label{tab:np09}\\
\toprule
\textbf{Necesidad} & \textbf{Frec.} \\
\midrule
\endfirsthead
\caption[]{Necesidad primaria: \emph{El robot navega en entornos domésticos.} (continuación).}\\
\toprule
\textbf{Necesidad} & \textbf{Frec.} \\
\midrule
\endhead
\midrule
\multicolumn{2}{r}{\small\emph{Continúa en la siguiente página}}\\
\endfoot
\bottomrule
\endlastfoot
\rowcolor{nprim} \textbf{El robot navega en entornos domésticos.} & \textbf{35} \\
\rowcolor{nsec} \textbf{El robot navega en entornos domésticos.} & \textbf{35} \\
El robot navega en entornos domésticos. & 8 \\
El robot navega evadiendo obstáculos. & 15 \\
El robot identifica obstáculos en el entorno. & 8 \\
El robot llega a su destino con exactitud. & 4 \\
\end{longtable}

\Needspace{9\baselineskip}
\begin{longtable}{@{}>{\raggedright\arraybackslash}p{0.85\textwidth} r@{}}
\caption{Necesidad primaria: \emph{El robot opera de forma continua por periodos prolongados.}}\label{tab:np10}\\
\toprule
\textbf{Necesidad} & \textbf{Frec.} \\
\midrule
\endfirsthead
\caption[]{Necesidad primaria: \emph{El robot opera de forma continua por periodos prolongados.} (continuación).}\\
\toprule
\textbf{Necesidad} & \textbf{Frec.} \\
\midrule
\endhead
\midrule
\multicolumn{2}{r}{\small\emph{Continúa en la siguiente página}}\\
\endfoot
\bottomrule
\endlastfoot
\rowcolor{nprim} \textbf{El robot opera de forma continua por periodos prolongados.} & \textbf{15} \\
\rowcolor{nsec} \textbf{El robot opera de forma continua por periodos prolongados.} & \textbf{15} \\
El robot optimiza el consumo de energía. & 5 \\
El robot opera de forma continua por periodos prolongados. & 9 \\
El robot requiere poco mantenimiento. & 1 \\
\end{longtable}

\Needspace{9\baselineskip}
\begin{longtable}{@{}>{\raggedright\arraybackslash}p{0.85\textwidth} r@{}}
\caption{Necesidad primaria: \emph{El robot tiene sensores confiables.}}\label{tab:np11}\\
\toprule
\textbf{Necesidad} & \textbf{Frec.} \\
\midrule
\endfirsthead
\caption[]{Necesidad primaria: \emph{El robot tiene sensores confiables.} (continuación).}\\
\toprule
\textbf{Necesidad} & \textbf{Frec.} \\
\midrule
\endhead
\midrule
\multicolumn{2}{r}{\small\emph{Continúa en la siguiente página}}\\
\endfoot
\bottomrule
\endlastfoot
\rowcolor{nprim} \textbf{El robot tiene sensores confiables.} & \textbf{18} \\
\rowcolor{nsec} \textbf{El robot tiene sensores confiables.} & \textbf{18} \\
El robot tiene sensores exactos y precisos. & 1 \\
El robot posiciona sus manipuladores con exactitud. & 10 \\
El robot opera bajo condiciones de iluminación variables. & 7 \\
\end{longtable}

\Needspace{8\baselineskip}
\begin{longtable}{@{}>{\raggedright\arraybackslash}p{0.85\textwidth} r@{}}
\caption{Necesidad primaria: \emph{El robot identifica la voz del usuario.}}\label{tab:np12}\\
\toprule
\textbf{Necesidad} & \textbf{Frec.} \\
\midrule
\endfirsthead
\caption[]{Necesidad primaria: \emph{El robot identifica la voz del usuario.} (continuación).}\\
\toprule
\textbf{Necesidad} & \textbf{Frec.} \\
\midrule
\endhead
\midrule
\multicolumn{2}{r}{\small\emph{Continúa en la siguiente página}}\\
\endfoot
\bottomrule
\endlastfoot
\rowcolor{nprim} \textbf{El robot identifica la voz del usuario.} & \textbf{5} \\
\rowcolor{nsec} \textbf{El robot identifica la voz del usuario.} & \textbf{5} \\
El robot registra sonido. & 3 \\
El robot identifica la voz del usuario. & 2 \\
\end{longtable}

\Needspace{7\baselineskip}
\begin{longtable}{@{}>{\raggedright\arraybackslash}p{0.85\textwidth} r@{}}
\caption{Necesidad primaria: \emph{El robot optimiza el tiempo de ejecución de sus tareas.}}\label{tab:np13}\\
\toprule
\textbf{Necesidad} & \textbf{Frec.} \\
\midrule
\endfirsthead
\caption[]{Necesidad primaria: \emph{El robot optimiza el tiempo de ejecución de sus tareas.} (continuación).}\\
\toprule
\textbf{Necesidad} & \textbf{Frec.} \\
\midrule
\endhead
\midrule
\multicolumn{2}{r}{\small\emph{Continúa en la siguiente página}}\\
\endfoot
\bottomrule
\endlastfoot
\rowcolor{nprim} \textbf{El robot optimiza el tiempo de ejecución de sus tareas.} & \textbf{1} \\
\rowcolor{nsec} \textbf{El robot optimiza el tiempo de ejecución de sus tareas.} & \textbf{1} \\
El robot optimiza el tiempo de ejecución de sus tareas. & 1 \\
\end{longtable}

\Needspace{7\baselineskip}
\begin{longtable}{@{}>{\raggedright\arraybackslash}p{0.85\textwidth} r@{}}
\caption{Necesidad primaria: \emph{El robot percibe profundidad.}}\label{tab:np14}\\
\toprule
\textbf{Necesidad} & \textbf{Frec.} \\
\midrule
\endfirsthead
\caption[]{Necesidad primaria: \emph{El robot percibe profundidad.} (continuación).}\\
\toprule
\textbf{Necesidad} & \textbf{Frec.} \\
\midrule
\endhead
\midrule
\multicolumn{2}{r}{\small\emph{Continúa en la siguiente página}}\\
\endfoot
\bottomrule
\endlastfoot
\rowcolor{nprim} \textbf{El robot percibe profundidad.} & \textbf{16} \\
\rowcolor{nsec} \textbf{El robot percibe profundidad.} & \textbf{16} \\
El robot percibe profundidad. & 16 \\
\end{longtable}
